\documentclass[11pt,letterpaper]{article}
\usepackage[margin=0.93in,top=0.85in,bottom=0.85in,headheight=14pt]{geometry}
\usepackage{amsmath,amssymb,amsthm,booktabs,array,microtype,xcolor,fancyhdr,graphicx,hyperref}
\microtypesetup{expansion=false}
\definecolor{ink}{HTML}{163748}
\definecolor{accent}{HTML}{14665D}
\hypersetup{colorlinks=true,linkcolor=accent,urlcolor=accent,citecolor=accent,
pdftitle={Optimal signings of square cycles and sharp Lubin-Tate jet recovery}}
\pagestyle{fancy}\fancyhf{}\fancyhead[L]{\small\sffamily\color{ink}Optimal signings and jet recovery}
\fancyhead[R]{\small\sffamily\color{ink}25 November 2026}\fancyfoot[C]{\thepage}
\setlength{\headsep}{16pt}\setlength{\parindent}{0pt}\setlength{\parskip}{5pt plus 1pt minus 1pt}
\setlength{\emergencystretch}{1.5em}
\newtheorem{theorem}{Theorem}[section]\newtheorem{lemma}[theorem]{Lemma}
\newtheorem{proposition}[theorem]{Proposition}\newtheorem{corollary}[theorem]{Corollary}
\newcommand{\TT}{\mathcal T}\newcommand{\RR}{\mathbb R}
\newcommand{\ZZ}{\mathbb Z}\newcommand{\eps}{\varepsilon}
\newcommand{\src}[2]{\href{https://github.com/lsernau/grh-lean/blob/72fead102f663d5c7a407964a13dc0cf9d977207/grh/src/#1}{#2}}
\newcommand{\code}[1]{{\small\texttt{#1}}}
\begin{document}
\thispagestyle{plain}
{\large\sffamily\color{accent}Research note \hfill 25 November 2026}\par
\vspace{0.5em}
{\fontsize{24.88}{30}\selectfont\bfseries\color{ink}Optimal signings of square cycles\par
and sharp Lubin--Tate jet recovery\par}
{\large\itshape Applications of methods from the Dirichlet GRH formalization}\par
\vspace{0.45em}\hrule\vspace{0.6em}
What is the smallest possible spectral radius of the adjacency matrix of the square of a cycle, when each edge can be signed independently?
For every $n=8m\ge32$ we find the exact minimum
\[
 \boxed{\ F(2\cos(\pi/m))\ },\qquad
 F(u)=\sqrt{4+\sqrt{8+u+\sqrt{26-3u}}}.
\]
Every minimizer has the same period-eight pattern of triangle signs up to rotation, and the product of signs around the ordinary cycle is $-1$. This determines the unrestricted optimum for the entire family in Suvagiya's Conjecture~28 \cite{Suvagiya}. The conjectured constant $r_*=F(2)$ is the second distinct spectral radius. A previous counterexample at $n=32$ attained a value below $2.79$ and left the exact optimum open \cite{Trure}. We obtain $2.78426979930755\ldots$ there.

The lower bound comes from an exact classification on 29 consecutive vertices. Of all possible length-27 triangle-sign words, only eight can have local norm below $1397/500$. They are the successive windows in one periodic word. Realizing both boundary conditions on that word then reduces the infinite family of finite problems to a single eight-dimensional determinant. The same argument resolves the three smaller sizes listed in Suvagiya's Remark~29 and gives an infinite-line rigidity theorem.

A second application gives the exact amount of information needed to reconstruct a bivariate Taylor polynomial from its substitutions along distinct Lubin--Tate branches. A compactness argument of Berger and Sprang gives the existence of a finite cutoff \cite{BS}. On a branch $b$ the least cutoff is
\[
                 1+\sum_{c\ne b}q^{v(b-c)}.
\]
Explicit vanishing products prove sharpness one branch at a time. We also determine optimal sets of directions and give a version over rings with nilpotents.

A further arithmetic consequence concerns Berger and Sprang's coefficient-depth invariant $s_0(n)$. A base-$p$ carry criterion determines when the first possible depth suffices for an infinite family of degrees. Along $n=(q-1)q^k$, $q=p^2$, the excess $s_0(n)-n$ is unbounded. Thus qualitative finiteness of every individual depth is compatible with an increasing cost of detecting certain degrees.

Both applications use ingredients encountered inside the GRH development. Their conclusions hold independently of the arithmetic zero-location theorems. The graph certificate uses exact rational and integer arithmetic. The interpolation result has a direct algebraic proof.

\newpage
\section{What transfers from the Lean proof}
The public project has a readable and a compact edition. We inspected the readable edition at commit \texttt{72fead102f66}, including its
\src{Statements.lean}{statement file} and \src{Main.lean}{five exported theorems}. It proves Mathlib's Riemann hypothesis, the critical-line formulation for $\zeta$, a right-half-plane theorem for all Dirichlet characters of positive modulus (with the principal pole exception), the critical-strip statement for all characters, and its primitive special case. The project reports kernel and independent export checks for Lean 4.35.0, pinned Mathlib, and the axioms \code{propext}, \code{Classical.choice}, and \code{Quot.sound}. This note uses source inspection. The new computations documented below were run separately \cite{Repo}.

\textbf{The arithmetic step.}
Short integral target vectors, a nearly full span and a lower height bound for a twisted determinant constrain the same index-corrected covolume. Its main term cancels between the two estimates. In the notation of
\src{New/Reduction/MainInequality.lean\#L87}{\code{star\_inequality}} the result is
\[
          m\,\mathrm{totalRemainder}\ \le\ f a (\log b)^5 .
\]
Here $m$ counts the successive lifts and lies between $\sqrt{x}$ and $2\sqrt{x}$. The choices of scales and the prime number theorem at a second, much larger cutoff give a one-sided estimate $R_{d,q}(x)\le x^{-1/2+\eps}$ eventually, for $18480\mid q$ and auxiliary primes $d\equiv q/2-1\pmod q$. The remainder smooths $\Lambda(n)(1-\chi_d(n))$ on the class $1\bmod q$ with the weight $1/n-1/x$, and subtracts $\log x/\varphi(q)$ and a constant. A quadratic twist avoids zeros in a prescribed rectangle. The surviving residues encode nonnegative multiplicities, and a one-sided Landau argument excludes them. See the \src{New/Reduction/OneSidedBound.lean}{upper bound} and the \src{New/Reduction/Zeros/Criterion.lean}{analytic criterion}. An extension to general Artin or automorphic $L$-functions would need new arithmetic estimates and twist control for its particular family.

\textbf{Two usable components.}
The proof of a block determinant estimate in \src{Lib/General/LinearAlgebra/HermitianMatrices.lean\#L50}{\code{HermitianMatrices}} restricts a positive form to a principal submatrix and takes a Schur complement. This classical operation supplies the inexpensive local obstruction in Section~\ref{sec:windows}. The roots-of-unity sum in \src{New/Derivatives/Vacuum/TightFrames.lean\#L388}{\code{isTightFrame\_pow}} is the orthogonality needed for the twisted Fourier decomposition. For the second application, \src{New/Derivatives/Sharpness/OneCurve.lean\#L45}{\code{OneCurve}} uses finite Taylor data and realizes an extremal coefficient on a monomial. A product with prescribed branch zeros gives an extremizer for the jet recovery problem. The assertions specific to our applications are the window classification, its unrestricted spectral consequences, and the sharp branch precisions.

\textbf{Earlier ProveIt work.}
The pinned incoming directory at \texttt{48efe71a726d} contains 41 archives. Its 248 TeX components were screened for overlap \cite{Incoming}. The $A_6$ and fixed-field reports construct detectors for nonabelian zeros. The sharp-torsion report optimizes valuation loss at division points, including Lubin--Tate torsion. Our variables in Section~\ref{sec:lubin} are formal endomorphism branches; the observable is an entire truncated substitution on each branch. We found no treatment of the signed-square-cycle extremum or this branch precision problem in the screened material.

\newpage
\section{The graph and its local obstruction}\label{sec:windows}
The vertices of $C_n(1,2)$ are $\ZZ/n\ZZ$, with a signed edge $a_i$ from $i$ to $i+1$ and $b_i$ from $i$ to $i+2$. Both are undirected. Its real symmetric signed adjacency $A$ has norm $\rho(A)=\max|\lambda(A)|$. Switching at vertices conjugates $A$ by a diagonal matrix with entries $\pm1$. Define
\[
       \tau_i=a_i a_{i+1}b_i,\qquad \alpha=\prod_{i\bmod n}a_i .
\]
These invariants completely specify a switching class on labeled vertices. Choose vertex signs successively to set $a_0,\ldots,a_{n-2}$ equal to $+1$. Then $a_{n-1}=\alpha$ and $b_i=\tau_i a_i a_{i+1}$. In particular, the last two distance-two edges must also be adjusted when $\alpha$ changes.

For a word $w=(t_0,\ldots,t_{k-1})$, let $H(w)$ be the $(k+2)$-vertex path with edges of length one and two, all length-one signs positive and length-two signs $t_i$. An interval of $k+2$ consecutive vertices is an induced copy when $n\ge k+4$: its endpoints are at wraparound distance at least three. A switch on that interval yields $H(w)$ for the corresponding consecutive triangle signs.
Put
\[
      T=1397/500=2.794,\qquad W=(+,+,-,+,-,-,+,-).
\]

\begin{lemma}[Finite obstruction]\label{lem:windows}
Exactly eight words $w$ of length $27$ have $\|H(w)\|<T$. They are the eight length-$27$ windows of $W$ repeated periodically. At length $19$ the only additional words are the two alternating windows.
\end{lemma}
\begin{proof}
A word passes precisely when both $TI-H(w)$ and $TI+H(w)$ are positive definite. Each leading principal submatrix must then be positive definite. Starting with the empty word, append each of the two signs to each surviving word. Testing only such extensions is exhaustive by induction.

Here is a rational recurrence for the test. After eliminating all but the last two vertices of $TI-H(w)$, write the remaining block as
$\left(\begin{smallmatrix}a&b\\b&c\end{smallmatrix}\right)$.
Appending sign $t$ adjoins a boundary column with entries $-t,-1$. Eliminating the vertex with pivot $a$ gives the three new entries
\[
     (a',b',c')=\bigl(c-b^2/a,\; -1+bt/a,\; T-1/a\bigr).
\]
Keep the extension if $a'>0$ and $a'c'>(b')^2$ for both shifted matrices. Alternating vertex signs transform $TI+H(w)$ into $TI-H(-w)$, so its second recurrence uses $-t$. Both start at $(T,-1,T)$. Evaluating this recurrence gives the two lists detailed on the following page.
\end{proof}


\newpage
\subsection*{Exact evaluation of the finite obstruction}
The recurrence of Lemma~\ref{lem:windows} gives the following counts. A prefix failing either strict inequality can have no admissible descendant.
\begin{center}\small
\begin{tabular}{r|rrrrrrrrrrrrrr}
\toprule
Length &0&1&2&3&4&5&6&7&8&9&10&11&12&13\\
Survivors &1&2&4&6&10&16&24&28&36&42&46&50&42&42\\
\midrule
Length &14&15&16&17&18&19&20&21&22&23&24&25&26&27\\
Survivors &44&38&32&22&18&10&10&10&10&10&10&10&10&8\\
\bottomrule
\end{tabular}
\end{center}
Counts alone do not identify the surviving words. The final set comparison is
\[
  \{w:\ |w|=27,\ \|H(w)\|<T\}
   =\bigl\{(W_{j+i\bmod 8})_{0\le i<27}:0\le j<8\bigr\},
\]
where the index on $W$ is read cyclically. Both sides are constructed explicitly by \code{verify\_graph.py}. It also compares the length-19 words with the indicated union. Its second calculation builds the integer matrices $1397I\pm500H(w)$ with the actual signs. For every surviving parent, positivity of the two new determinants is equivalent to both leading-minor tests, since all earlier leading minors are already positive. It recomputes those determinants by fraction-free integer elimination. A third implementation in \code{verify\_by\_permutations.py} sums signed permutation products and obtains the same eight words. The certificate takes no eigenvalue tolerance or external input.

\begin{proposition}[Rigidity below the cutoff]\label{prop:rigidity}
If $n\ge31$ and $\rho(A)<T$, then $8\mid n$ and the triangle signs form a rotation of repeated $W$. When $n\ge23$ the length-19 test forces either this pattern or alternation.
\end{proposition}
\begin{proof}
Compression of a real symmetric matrix to a coordinate subspace has norm at most the original norm, by the Rayleigh quotient. Every length-27 window must therefore be in the finite list if $n\ge31$. Its first 26 signs determine its phase, since the eight rotations already have distinct first eight signs. Overlap forces the next letter at every step. The primitive period is eight, so closing the word around the cycle requires $8\mid n$. For the length-19 list, the ten prefixes of length 18 are distinct. The unique continuations consist of one cycle of length eight and one of length two.
\end{proof}

At $n=32$ there are $2^{33}$ labeled switching classes before the test. The proposition leaves sixteen possibilities below the threshold, one rotation for each choice of $\alpha$. Their spectra are computed next. The finite bound applies at every larger length with exactly the same proof.

\newpage
\section{Exact optimum from the two boundary conditions}
Normalize the length-one edges on the covering line. For repeated $W$ the adjacency acts by
\[
 (Af)_i=f_{i-1}+f_{i+1}+W_{i-2} f_{i-2}+W_i f_{i+2}.
\]
The condition after one circuit is $f_{i+8m}=\alpha f_i$. Diagonalize translation by eight. On the block indexed by $z^m=\alpha$, $f_{i+8}=zf_i$. The $m$ blocks are mutually orthogonal because $\sum_{j=0}^{m-1}(z/ z')^j=0$ for $z\ne z'$. Their dimensions add to $8m$.

Suvagiya's equation (3) gives the characteristic polynomial for the eight-vertex block \cite{Suvagiya}. The same symbol applies to either holonomy. An explicit check in the archive builds the symbol from the signed edges, treating $z$ as an indeterminate, and expands its determinant. With $u=z+z^{-1}$ and $h=x^2-4$ the polynomial simplifies to
\[
          h^4-(16+2u)h^2+u^2+19u+38.
\]
Indeed the two values of $h^2$ are $8+u\pm\sqrt{26-3u}$. The greatest possible $x^2$ is therefore $4+\sqrt{8+u+\sqrt{26-3u}}$ and the block norm is $F(u)$. The function increases strictly on $[-2,2]$, since
\[
 \frac{d}{du}\big(8+u+\sqrt{26-3u}\big)
 =1-\frac{3}{2\sqrt{26-3u}}>0.
\]
In the periodic case $z=1$ attains $u=2$. In the antiperiodic case the largest allowed value is $2\cos(\pi/m)$.

\begin{theorem}[Unrestricted minimum and all minimizers]\label{thm:main}
For every $m\ge4$ the minimum of $\rho(A)$ over all signings of $C_{8m}(1,2)$ is $F(2\cos(\pi/m))$. Precisely eight labeled switching classes attain it, with triangle word a rotation of $W$ and $\alpha=-1$. The next distinct value is $r_*=F(2)$, attained by eight classes with the same rotations and $\alpha=+1$. All remaining classes have $\rho(A)\ge1397/500$.
\end{theorem}
\begin{proof}
The two candidate radii are strictly ordered by monotonicity. Both are below $T$: $\sqrt{20}<4.473$, hence $F(2)^2=4+\sqrt{10+\sqrt{20}}<7.805<T^2$. Proposition~\ref{prop:rigidity} excludes every other word from this interval. Completeness of the switching invariants gives the exact counts.
\end{proof}
For $n=32$, $u=\sqrt2$ at the winning phase. The optimizer is precisely the switching orbit, up to rotation, of the earlier counterexample \cite{Trure}. Its optimality now follows for arbitrary choices on all 64 edges. As $m$ increases, the minimum increases strictly to $r_*$, with
$ r_*-F(2\cos(\pi/m))=\pi^2 F'(2)m^{-2}+O(m^{-4})$.

\newpage
\section{Consequences for other lengths and graph zeta functions}
\textbf{The first even length at which alternation loses.}
For every even $8\le n\le30$, the exact unrestricted minimum is
\[
           G(n)=2\sqrt{\cos^2(\pi/n)+\cos^2(2\pi/n)} .
\]
Suvagiya's Proposition~25 computes this value at negative holonomy in the alternating family and gives $2\sqrt2$ at positive holonomy. His Remark~29 records heuristic searches at $n=26,28,30$ \cite{Suvagiya}. The corresponding minima are $2.777192$, $2.784204$ and $2.789872$ to six decimals.

For $n\ge16$, a shorter cyclic version of the obstruction suffices. Regard each of the 42 surviving length-12 words at $T$ as an arrow from its 11-letter prefix to its 11-letter suffix. Delete any arrow whose initial vertex has no incoming arrow or whose final vertex has no outgoing arrow. No arrow on a closed walk can be removed. The counts after successive rounds are $42,20,10,10$. The ten remaining arrows form the period-eight and alternating cycles. This exact computation is included in \code{verify\_by\_permutations.py}. Each length-12 word occupies 14 vertices, giving the claimed range of $n$.

The increasing function $G(n)$ is below $T$ for $n\le30$: put $v=3.14/30$ and $C(t)=1-t^2/2+t^4/24$. Then $4(C(v)^2+C(2v)^2)<T^2$ by integer arithmetic, and $C(v),C(2v)$ bound the cosines at $n=30$ from above. Period eight is compatible only with $n=16,24$ in the range $16\le n\le30$. At 16, $G(16)^2=4+\sqrt{2+\sqrt2}+\sqrt2<7.3$, while $F(0)^2>7.6$. At 24, $G(24)^2=4+(\sqrt6+\sqrt2)/2+\sqrt3<7.67$, while $F(1)^2>7.7$. This proves the assertion throughout that range. The smaller cases $8,10,12,14$ have direct finite certificates in \code{verify\_small\_cycles.py}.

\begin{center}\includegraphics[width=.99\linewidth]{cycle_optima.pdf}\end{center}
{\small The left panel shows the alternating optimum at every even length in its indicated range. The right panel uses a finer vertical scale. The winning negative-holonomy radius tends to the positive-holonomy radius from below.}\par
\newpage
\subsection*{Infinite graphs and nontrivial Ihara poles}
\textbf{A local gap and the line.}
A stronger test at $T'=281/100$ is short enough to reproduce in the independent verifier. There are 72 surviving length-43 words. Their central nine-letter subwords (indices 17 through 25 from zero) are exactly the eight windows of $W$. Thus for $n\ge47$ any signing whose triangle word differs from the period-eight family has norm at least $2.81$. The same statement holds for signings of the graph on $\ZZ$ joining integers at distance one or two: every finite compression is an operator-norm lower bound, and every position is central in a length-43 window. On the infinite line all length-one signs can be switched to $+1$. Fourier transformation by period-eight cells gives the norm of the surviving operator as $\max_{-2\le u\le2}F(u)=r_*$. This proves an isolated minimum on the line, uniquely up to switching and translation.

\textbf{An extremal statement about poles.}
For a signed four-regular graph, the Bass determinant formula for the signed Ihara function gives the nontrivial factor of its reciprocal as
\[
         Q_A(v)=\det(I-vA+3v^2 I)
                  =\prod_{\lambda\in\mathrm{spec}(A)}(1-\lambda v+3v^2).
\]
The universal factor $(1-v^2)^{|E|-|V|}$ is omitted in $Q_A$. This form also follows by taking the ordinary Bass identities for the double cover and base graph and dividing them: adjacency on the cover splits into the base and signed matrices \cite{Suvagiya,Bass}. When $|\lambda|\le2\sqrt3$ the two roots of its factor are $3^{-1/2}e^{\pm i\theta}$, where $2\sqrt3\cos\theta=\lambda$. For $|\lambda|>2\sqrt3$ there is a root on the real axis. Thus, among all signings of $C_{8m}(1,2)$, the largest possible minimum angular distance of the roots of $Q_A$ to the real axis is
\[
           \arccos\!\left(\frac{F(2\cos(\pi/m))}{2\sqrt3}\right).
\]
Theorem~\ref{thm:main} characterizes every maximizing signing. At $n=32$ the angle is about $36.51^\circ$. The figure shows the roots without multiplicities. The signed Ihara function itself has poles at these points.
\begin{center}\includegraphics[width=.78\linewidth]{ihara_poles.pdf}\end{center}

\newpage
\input{lubin_tate}
\newpage
\input{first_depth}
\section{Further research questions}
\begin{enumerate}
\item What are the exact optima at large lengths incompatible with period eight? The threshold $2.81$ excludes a whole interval at lengths $n\ge47$ with $8\nmid n$. A larger cutoff could reveal the allowed local departures from periodicity and their effect on the spectrum.
\item Can the 29-vertex test be replaced by a short human-checkable family of inequalities for forbidden sign patterns? Such a proof would explain the period eight and could extend to $C_n(1,k)$ with $k>2$.
\item For recovery modulo $\pi^r$, which branch configurations attain the pole-order bound in Proposition~\ref{prop:integralbranch}? With more than two branches, products of inverse differences can exhibit cancellation in their highest pole terms.
\item What is the true growth of $s_0((q-1)q^k)-(q-1)q^k$? The obstruction in Section~\ref{sec:first-depth} supplies a growing lower bound. A matching upper bound would require the explicit construction of trace-kernel rows at a controlled depth.
\end{enumerate}

\textbf{Verification and priority.}
The graph proof has an exact computational lemma and an analytic argument valid at every length in the stated family. The jet proof is algebraic. The accompanying scripts independently exercise the finite calculations. Neither new result has been formalized in Lean or externally refereed here. The literature comparison concerns the concrete prior claims identified in the text. Publication priority is unverified.

\begin{thebibliography}{9}\small
\bibitem{Repo} The \texttt{lsernau/grh-lean} contributors. \emph{GRH for Dirichlet $L$-functions, in Lean 4}. Readable edition, commit \texttt{72fead102f663d5c7a407964a13dc0cf9d977207}. \href{https://github.com/lsernau/grh-lean/tree/72fead102f663d5c7a407964a13dc0cf9d977207/grh}{Source, guide and reported machine checks}. Original method attributed there to L. Alp\"oge with Claude.
\bibitem{Incoming} V. Reshetnikov, \emph{ProveIt}, \href{https://github.com/VladimirReshetnikov/ProveIt/tree/48efe71a726dddc5a505534edaa4c0fd190416b9/docs/incoming}{incoming research reports}, commit \texttt{48efe71a726d}. File names and blob hashes accompany this note.
\bibitem{Suvagiya} V. Suvagiya. \emph{Parity families and signed spectra: kernel averaging, near-Ramanujan bounds, and exact circulant models}. \href{https://arxiv.org/html/2607.17343v2}{arXiv:2607.17343v2}, esp. Propositions 23, 25, Theorem 26, Remarks 27, 29 and Conjecture 28 (2026).
\bibitem{Trure} trureturing. \emph{Suvagiya Conjecture 28 refutation}. \href{https://www.trureturing.org/research/suvagiya-conjecture28-refutation/}{Public problem dossier and 32-vertex witness} (2026).
\bibitem{Bass} H. Bass. The Ihara--Selberg zeta function of a tree lattice. \emph{Int. J. Math.} 3 (1992), 717--797. \href{https://doi.org/10.1142/S0129167X92000357}{Publisher}.
\bibitem{BS} L. Berger and J. Sprang. \emph{Integer-valued polynomials and $p$-adic Fourier theory}. \href{https://arxiv.org/abs/2502.18053}{arXiv:2502.18053}, Lemma 4.2, Proposition 4.4 and Theorem 4.6 (2025).
\bibitem{BSS} L. Berger, J. Sprang and R. Steingart. \emph{Bounded functions on the character variety for quadratic extensions of $\mathbb Q_p$}. \href{https://arxiv.org/abs/2611.21796}{arXiv:2611.21796} (2026).
\bibitem{Colovic} U. A. Colovi\'c. \emph{On the ring of bounded functions on the character variety}. \href{https://arxiv.org/abs/2611.25162}{arXiv:2611.25162} (2026).
\bibitem{Bhargava} M. Bhargava. $P$-orderings and polynomial functions on arbitrary subsets of Dedekind rings. \emph{J. reine angew. Math.} 490 (1997), 101--127. \href{https://doi.org/10.1515/crll.1997.490.101}{Publisher}.
\end{thebibliography}
\end{document}
